% !TeX root = ../main.tex

\begin{abstract}

隨著 5G-Advanced 網路部署密度與大規模多輸入多輸出（Massive MIMO，如 64T64R）技術的指數級增長，基地台的高耗電問題已成為電信營運商降低營運成本與達成淨零碳排目標的首要挑戰。在 3GPP 研議的網路節能（Network Energy Saving, NES）技術中，空間域適應（Spatial Domain Adaptation, SDA）透過動態增減射頻鏈與天線數量來降低功耗，被視為極具潛力的解決方案。然而，傳統基於負載（Load-based）或靜態規則的天線調適策略，在面對實際網路中的突發性流量（Burst Traffic）時往往反應遲鈍，難以在「節能效益」與「服務品質（QoS）」的兩難中取得最佳平衡，容易導致嚴重的佇列堆積與網路效能震盪。

為解決此研究缺口，本論文提出一套「流量感知的動態權重李雅普諾夫最佳化演算法」（Traffic-Aware Dynamic V-Control Lyapunov Optimization）。本研究跳脫傳統的貪婪啟發式規則，引入李雅普諾夫最佳化理論作為決策大腦，利用漂移加懲罰（Drift-Plus-Penalty, DPP）數學框架，將 MAC 層的佇列長度與 PHY 層的基地台功耗進行跨時間尺度的聯合最佳化。此外，本演算法創新地提出了動態權重（Dynamic $V$）控制機制，能根據即時流量變化自動調整系統對「省電」的執著程度，並結合 MAC-PHY 跨層映射（Bits to RU），在不依賴未來先驗知識的條件下，提供具備數學防護網的即時天線收放決策。

在系統驗證方面，本研究採用兩階段驗證方法論：第一階段透過 MATLAB 建立演算法沙盒進行理論收斂性與參數調優；第二階段則將最佳化之演算法模組移植至符合 3GPP TR 38.864 標準、基於 7-site 21-cell 網格拓撲之高保真度系統級模擬器（SLS）中進行全網實戰驗證。實驗結果表明，本研究提案在複雜的多細胞同頻干擾現網環境中，成功斬獲高達 \textbf{30.3\% 的長期平均節能率}。雖然受限於現網波束拓寬之同頻干擾外溢與實體層 CSI 量測時滯，使系統在\textbf{用戶感知吞吐量（UPT）與封包排隊延遲指標上承受了部分可控之性能折衷}，但所有性能指標均被\textbf{嚴格限制於李雅普諾夫強穩定性的理論安全邊界內}，且有效消除了傳統 Heuristic 機制常見的天線高頻乒乓開關效應。本方案以極低且理論可控的 QoS 代價換取了深度綠能紅利，極具產業落地價值，完美契合未來 5G-Advanced 與 6G 網路針對低碳、省電背景切片技術之部署願景。

\end{abstract}

\clearpage

\begin{abstract*}

With the exponential growth of 5G-Advanced network deployment density and Massive MIMO (e.g., 64T64R) technologies, the high power consumption of base stations (BSs) has become a primary challenge for telecom operators seeking to reduce operational expenditure (OPEX) and achieve net-zero carbon emission targets. Among the Network Energy Saving (NES) techniques investigated by 3GPP, Spatial Domain Adaptation (SDA) is considered a highly promising solution. However, traditional load-based or static rule-driven antenna adaptation strategies respond sluggishly to burst traffic, struggling to achieve an optimal balance in the trade-off between energy-saving efficiency and service quality (QoS). 

To bridge this research gap, this paper proposes a Traffic-Aware Dynamic V-Control Lyapunov Optimization algorithm. Utilizing the Drift-Plus-Penalty (DPP) mathematical framework, the MAC-layer queue length (delay) and the PHY-layer BS power consumption are jointly optimized. Furthermore, the algorithm innovatively incorporates a dynamic weight control mechanism that automatically adjusts the system's emphasis on energy saving based on real-time traffic fluctuations. Combined with a cross-layer MAC-PHY mapping, the proposed scheme provides real-time antenna adaptation decisions backed by mathematical stability guarantees, without relying on future traffic a priori knowledge. 

For system validation, a rigorous two-stage evaluation methodology is conducted, spanning a logical MATLAB sandbox for parameter tuning and a high-fidelity C++ System-Level Simulator (SLS) aligned with the 3GPP TR 38.864 standard under a 21-cell wraparound network topology. Simulation results demonstrate that the proposed algorithm precisely captures the mathematically bounded trade-off between energy efficiency and delay. Compared to traditional SDA baselines, the proposed design achieves an outstanding \textbf{Energy Saving Gain (ESG) of up to 30.3\%}. Although co-channel beam-widening interference and optimistic channel state information (CSI) estimation introduce a manageable degradation in queueing delay and User-Perceived Throughput (UPT), \textbf{all performance metrics are strictly confined within the theoretically-proven Lyapunov strong stability region}. This demonstrates the practical value of trading controllable, non-divergent QoS degradation for maximum green energy dividends, making the proposed algorithm highly suitable for sustainable, deep-green background slices in 5G-Advanced and 6G networks.

\end{abstract*}